\documentclass[../notes.tex]{subfiles}
\graphicspath{
    {"../figures"}
}

\begin{document}

\subsection{Variation of Parameters}

\begin{theorem}{Variation of Parameters}{varofparams}
				Given a \nth{2} Order ODE $y'' + p(t) y' + q(t) y = g(t)$, the particular solution of the ODE can be written in the form of
				\[
				y_p = u_1 y_1 + u_2 y_2
				.\] 
				where $y_1$ and $y_2$ are the complementary function solutions to the homogeneous version of the ODE and $u_1$ and $u_2$ are functions of $t$ that follow the following conditions:
				\begin{align*}
								 u_1' y_1 + u_2' y_2 &= 0 \\
								 u_1' y_1' + u_2' y_2' &= g(t)
				.\end{align*}
\end{theorem}

The benefit to \namethmref{th:varofparams} is that derivatives of $g(t)$ aren't the limiting factor, but rather the ability to find the functions $u_1$ and $u_2$ is the limiting factor. This means \thmref{th:varofparams} is able to find particular solutions in cases where the method of undetermined coefficients is unable to do so.

\centering
\begin{tikzpicture}[scale=8,every node/.style={transform shape}]
    \node [rotate=30] {Midterm};    
\end{tikzpicture}

\end{document}
